\chapter{Análisis y Especificación de Requisitos}
\label{cap:analisisRequisitos}

Este capítulo detalla a quién va dirigida la aplicación, qué roles intervienen y qué funcionalidades y restricciones debe cubrir la solución.

\section{Personas}

\subsection{Persona 1: Klaudia, Escaladora Constante}

Perfil de referencia para representar a una usuaria frecuente de rocódromo que necesita registrar su progreso y mantener la motivación.

\section{Roles}

\subsection{Escalador}

Rol principal orientado al uso diario de la aplicación para registrar actividad, consultar estadísticas y participar en la gamificación.

\subsection{Gestor de Rocódromo}

Rol encargado de administrar información del centro, rutas y elementos de seguimiento asociados al rocódromo.

\section{Especificación de funcionalidades (Historias de Usuario / Casos de Uso)}

Se agrupan aquí las funcionalidades que la aplicación debe ofrecer, redactadas como historias de usuario o como casos de uso según convenga al nivel de detalle.

\section{Requisitos no funcionales}

Se incluyen las expectativas relativas a rendimiento, seguridad, mantenibilidad, usabilidad y compatibilidad.

\section{Alcance de las versiones}

\subsection{Requisitos del prototipo (v.0.1.0)}

Se describen las capacidades mínimas del primer prototipo funcional.

\subsection{Requisito del MVP (v.1.0.0)}

Se detallan las capacidades que debe alcanzar la primera versión estable y utilizable del producto.